\begin{exercise}
Do research on the Pigeonhole Principle.  State it, and give two examples of how it can be used to solve an counting problem.  Are there other historical names for this counting principle?
\end{exercise}

\begin{exercise} 
Many ways are there for 8 different men and 5 different women to stand in a line so that no two women stand next to each other? What is the answer if there are 10 men and 6 women?  What is the answer if there are $m$ men and $w$ women with $m\geq w$?  How does the answer change if you cannot tell the men apart and you cannot tell the women apart, for example if the people were replace by mannequins.
\end{exercise}

\begin{exercise} 
A {\bf diagonal} of a polygon is a line joining two nonadjacent vertices.  How may diagonals does a hexagon have? What about a octogon? What about a $n$-gon?
\end{exercise}

\begin{exercise}
Suppose you have 10 points in the plane with no three collinear.  How many different line segments are formed by joining pairs of  of these points? How many different triangles are formed by these line segments?  If A is one of the 10 points, how many of the triangles formed by the line segments have A as a vertex?
\end{exercise}

\begin{exercise}
Find the definition of a multinomial coefficient.  How is it related to binomial coefficients?  What can a multinomial coefficient be used to count?  Give examples of counting with multinomial coefficients.
\end{exercise}

\begin{exercise}
How many different strings can be made from the letters in the work PEPPERCORN when all the letters are used?  How may of the strings start and end with the letter P?  How many strings have 3 consecutive Ps?
\end{exercise}

\begin{exercise}
How many length 9 bit strings are there which:
Start with the sub-string 101? \\
Have weight 5 (i.e., contain exactly five 1’s) and start with the sub-string 101? \\
Have weight 5 and either start with 101 or end with 11 (or both)? \\
\end{exercise}

\begin{exercise}
How many bit strings of length 10 are there with exactly 4 ones?  How many bit strings are there of length 13 are there with exactly 7 ones?  How many bit strings of length $n$ are there with exactly $k$ ones?
\end{exercise}

\begin{proofp}
Use mathematical induction to show that 21 divides $4^{n+1}+5^{2n-1}$ for any positive integer $n$
\end{proofp}

\begin{proofp}
Geometric sequence: Use induction to prove that for any real numbers $a$ and $r$, $$\sum_{i=0}^{n} a\cdot r^i = \frac{a(1-r^{n+1})}{1-r}.$$
\end{proofp}

\begin{proofp}
Arithmetic sequence: Use induction to prove that for any real numbers $a$ and $r$, $$\sum_{i=0}^{n} (a+ i\cdot r) = \frac{n+1}{2}(2a+n\cdot r).$$
\end{proofp}

\begin{proofp}
Use induction to prove that $8^n-3^n$ is divisible by 5 for all integers $n \ge 1$.  Hint:  This is equivalent to showing $8^n-3^n=5k$ for some integer $k$.
\end{proofp}


\begin{proofp}
Use induction to prove that $n! \ge 2^{n-1}$ for for all integers $n \ge 1$.
\end{proofp}

\begin{proofp}
Use induction to prove that $1+3+3^2+\cdots +3^n=(3^{n+1}-1)/2$.
\end{proofp}

\begin{proofp}
Prove by induction that $3^n+7^n-2$ is divisible by 8 for all positive integers $n$. Hint:  This is equivalent to showing $3^n+7^n-2=8k$ for some integer $k$.
\end{proofp}

\begin{proofp}
Prove $2^n>n^3$ for $n \ge 10$.
\end{proofp}

\begin{proofp}
Prove $n^2 \geq2n+3$ for $n \ge 3$.
\end{proofp}

\begin{proofp}
Prove $2^n>n^2$ for $n \ge 5$.
\end{proofp}

\begin{proofp}
By experimenting with small values of $n$, guess a formula for the given sum,
$$\frac{1}{1\cdot2}+\frac{1}{2\cdot3}+\frac{1}{3\cdot4}+\cdots+\frac{1}{n(n+1)};$$
then use induction to prove your formula.
\end{proofp}

\begin{proofp}$\bigstar$
 Bernoulli's inequality: Prove that if $h>-1$, then $1+nh \leq  (h+1)^n$ for all non-negative integers $n$. 
\end{proofp}

\begin{proofp}$\bigstar$
Use induction to prove that $n^2+n+1$ is odd for every integer $n$. 
\end{proofp}

%\begin{proofp}
%Prove that for $n \ge 0$, $2^n>n$.
%\end{proofp}

%\begin{proofp}
%Prove that for $n \ge 4$, $n!>2^n$.
%\end{proofp}

\begin{proofp}$\bigstar$
Prove by induction that $$\frac{1 \cdot 3 \cdot 5 \cdots (2n-1)}{2\cdot 4 \cdot 6 \cdots (2n)} \leq \frac{1}{\sqrt{n+1}}$$ for $n$ a positive integer.
\end{proofp}


\begin{proofp}$\bigstar$
Prove that for $n \ge 2$, the maximum number of points of intersection of $n$ distinct lines in the plane is $n(n-1)/2$
\end{proofp}


